\documentclass[book.tex]{subfiles}
\begin{document}

\chapter{Smooth Overlap of Atomic Positions}

\label{chap:soap}
\noindent\rule{11cm}{0.4pt} \\
\textbf{Prerequisites:} \autoref{chap:dft} \\
\textbf{Difficulty Level:} ** \\
\noindent\rule{11cm}{0.4pt}
\vspace{5mm}

The smooth overlap of atomic positions proposes a kernel method which uses a similarity notion between atomic neighborhoods to construct estimates of energies.

\section{Mathematical Framework}

The similarity of two atomic environments is given by the formula

\begin{align}
S(\rho, \rho') &= \int \rho(r) \rho'(r) dr
\end{align}
We can transform this definition into a rotationally invariant kernel by integrating over all choices of rotation matrix
\begin{align}
k(\rho, \rho') &= \int \left | S(\rho, R \rho') \right |^n  dR \\
&= \int \left | \int \rho(r) \rho'(R r)\right |^n  dR
\end{align}

The atomic densities are constructed as a sum of atom centered Gaussians
\begin{align}
\rho(r) &= \sum_i \exp(-\alpha |r - r_i|^2)
\end{align}
We can use the fact that the spherical harmonics provide an orthogonal basis to expand
\begin{align}
\exp(-\alpha|r - r_i|^2) &= \sum_{lm} c^i_{lm} Y_{lm}(r)
\end{align}
The overlap kernel can then be computed directly in terms of Wigner-D matrices using the fact that these matrices act on spherical harmonics in a structured way.

The full SOAP kernel is typically normalized as well

\begin{align}
K(\rho, \rho') &= \left (\frac{k(\rho,\rho')}{\sqrt{k(\rho, \rho)k(\rho', \rho')}} \right )^\zeta
\end{align}

This kernel method turns out to correspond directly to some of the featurization methods we have seen previously. However, the choice of a general kernel reduces some of the overhead associated with this system.

\section{Exercises}

\begin{enumerate}
    \item 
\end{enumerate}
\end{document}


